% ======================================================================
% Chapter 6: Radiative Corrections: Introduction
% (book pages 175-210; PDF pages 198-233)
% ======================================================================
\chapter{Radiative Corrections: Introduction}

Now that we have acquired some experience at performing QED calculations,
let us move on to some more complicated problems.  Chapter~5 dealt only with
tree-level processes, that is, with diagrams that contain no loops.  But all
such processes receive higher-order contributions, known as \emph{radiative
corrections}, from diagrams that do contain loops.  Another source of
radiative corrections in QED is \emph{bremsstrahlung}, the emission of extra
final-state photons during a reaction.  In this chapter we will investigate
both types of radiative corrections, and find that it is inconsistent to
include one without also including the other.

Throughout this chapter, in order to illustrate these ideas in the simplest
possible context, we will consider the process of electron scattering from
another, very heavy, particle.  We analyzed this process at tree level in
Section~5.4 and Problem~5.1.  At the next order in perturbation theory, we
encounter the following four diagrams:
\begin{equation}
  \text{(vertex correction)} + \text{(two external leg corrections)}
    + \text{(vacuum polarization)} .
  \label{eq:6.1}
\end{equation}
The order-$\alpha$ correction to the cross section comes from the interference
term between these diagrams and the tree-level diagram.  There are six
additional one-loop diagrams involving the heavy particle in the loop, but
they can be neglected in the limit where that particle is much heavier than
the electron, since the mass appears in the denominator of the propagator.
(Physically, the heavy particle accelerates less, and therefore radiates less,
during the collision.)

Of the four diagrams in \eqref{eq:6.1}, the first (known as the
\emph{vertex correction}) is the most intricate and gives the largest variety
of new effects.  For example, it gives rise to an anomalous magnetic moment
for the electron, which we will compute in Section~6.3.

The next two diagrams of \eqref{eq:6.1} are \emph{external leg corrections}.
We will neglect them in this chapter because they are not amputated, as
required by our formula Eq.~(4.90) for $S$-matrix elements.  We will
discuss these diagrams in more detail when we prove that formula in
Section~7.2.

The final diagram of \eqref{eq:6.1} is called the \emph{vacuum
polarization}.  Since it requires more computational machinery than the
others, we will not evaluate this diagram until Section~7.5.

Our study of these corrections will be complicated by the fact that they are
ill-defined.  Each diagram of \eqref{eq:6.1} involves an integration over the
undetermined loop momentum, and in each case the integral is divergent in the
$k\to\infty$ or \emph{ultraviolet} region.  Fortunately, the infinite parts of
these integrals will always cancel out of expressions for observable quantities
such as cross sections.

The first three diagrams of \eqref{eq:6.1} also contain \emph{infrared
divergences}: infinities coming from the $k\to 0$ end of the loop-momentum
integrals.  We will see in Section~6.4 that these divergences are canceled
when we also include the following \emph{bremsstrahlung} diagrams:
\begin{equation}
  \text{photon emission by initial electron} + \text{photon emission by
    final electron} .
  \label{eq:6.2}
\end{equation}
These diagrams are divergent in the limit where the energy of the radiated
photon tends to zero.  In this limit, the photon cannot be observed by any
physical detector, so it makes sense to add the cross section for producing
these low-energy photons to the cross section for scattering without
radiation.  The bremsstrahlung diagrams are thus an essential part of the
radiative correction, in this and any other QED process.

Our main goals in the present chapter are to understand bremsstrahlung of
low-energy photons, the vertex correction diagram, and the cancellation of
infrared divergences between these two types of radiative corrections.

\section{Soft Bremsstrahlung}

Let us begin our study of radiative corrections by analyzing the
bremsstrahlung process.  In this section we will first do a classical
computation of the intensity of the low-frequency bremsstrahlung radiation
when an electron undergoes a sudden acceleration.  We will then compute a
closely related quantity in quantum field theory: the cross section for
emission of one very soft photon, given by diagrams \eqref{eq:6.2}.  We would
like to understand how the classical result arises as a limiting case of the
quantum result.

\subsection{Classical Computation}

Suppose that a classical electron receives a sudden kick at time $t=0$ and
position $\mathbf{x}=0$, causing its 4-momentum to change from $p$ to $p'$.
(An infinitely sudden change of momentum is of course an unrealistic
idealization.  The precise form of the trajectory during the acceleration
does not affect the low-frequency radiation, however.  Our calculation will be
valid for radiation with a frequency less than the reciprocal of the
scattering time.)

We can find the radiation field by writing down the current of this electron,
and considering that current as a source for Maxwell's equations.

What is the current density of such a particle?  For a charged particle at
rest at $\mathbf{x}=0$, the current would be
\begin{equation}
  j^\mu(\mathbf{x}) = e\,\delta^{(3)}(\mathbf{x})\,
    \begin{pmatrix}1\\ 0\end{pmatrix}^{\mu}
  = e\int dt\ \begin{pmatrix}1\\ 0\end{pmatrix}^{\mu}\,
    \delta^{(3)}\bigl( \mathbf{x} - \mathbf{y}(t) \bigr)
  \qquad\text{with } y^\mu(t) = (t, 0)^\mu .
\end{equation}
From this we can guess the current for an arbitrary trajectory $y^\mu(\tau)$:
\begin{equation}
  j^\mu(x) = e\int d\tau\ \frac{dy^\mu}{d\tau}\,
    \delta^{(4)}\bigl( x - y(\tau) \bigr).
  \label{eq:6.3}
\end{equation}
Note that this expression is independent of the precise way in which the curve
$y^\mu(\tau)$ is parametrized: Changing variables from $\tau$ to
$a(\tau)$ gives a factor of $d\tau/da$ in the integration measure, which
combines with $dy^\mu/d\tau$ via the chain rule to give $dy^\mu/da$.  We can
also prove from \eqref{eq:6.3} that the current is automatically conserved:
For any ``test function'' $f(x)$ that falls off at infinity, we have
\begin{equation}
  \begin{aligned}
    \int d^4x\ f(x)\,\partial_\mu j^\mu(x)
    &= \int d^4x\ f(x)\,e\int d\tau\
      \frac{dy^\mu}{d\tau}\ \frac{\partial}{\partial x^\mu}
      \delta^{(4)}(x-y(\tau))\\
    &= -e\int d\tau\ \frac{dy^\mu}{d\tau}\ \frac{\partial f}{\partial x^\mu}
      \Big|_{x=y(\tau)}
    = -e\int d\tau\ \frac{d}{d\tau}f\bigl( y(\tau) \bigr)\\
    &= -e\Bigl[ f\bigl( y(\tau) \bigr) \Bigr]_{-\infty}^{\infty} = 0 .
  \end{aligned}
\end{equation}
For our process the trajectory is
\begin{equation}
  y^\mu(\tau) = \begin{cases}
    (p^\mu/m)\,\tau & \text{for } \tau < 0;\\
    (p'^{\mu}/m)\,\tau & \text{for } \tau > 0 .
  \end{cases}
\end{equation}
Thus the current can be written
\begin{equation}
  j^\mu(x) = e\int_0^{\infty} d\tau\ \frac{p'^{\mu}}{m}\,
    \delta^{(4)}\Bigl( x - \frac{p'}{m}\tau \Bigr)
    + e\int_{-\infty}^{0} d\tau\ \frac{p^{\mu}}{m}\,
      \delta^{(4)}\Bigl( x - \frac{p}{m}\tau \Bigr).
\end{equation}
In a moment we will need to know the Fourier transform of this function.
Inserting factors of $e^{-\epsilon\tau}$ and $e^{\epsilon\tau}$ to make the
integrals converge, we have
\begin{equation}
  \begin{aligned}
    \tilde{j}^\mu(k) &= \int d^4x\ e^{ik\cdot x}\,j^\mu(x)\\
    &= e\int_0^{\infty} d\tau\ \frac{p'^{\mu}}{m}\,
      e^{i(k\cdot p'/m+i\epsilon)\tau}
      + e\int_{-\infty}^{0} d\tau\ \frac{p^{\mu}}{m}\,
        e^{i(k\cdot p/m-i\epsilon)\tau}\\
    &= ie\Bigl( \frac{p'^{\mu}}{k\cdot p' + i\epsilon}
      - \frac{p^{\mu}}{k\cdot p - i\epsilon} \Bigr).
  \end{aligned}
  \label{eq:6.4}
\end{equation}
We are now ready to solve Maxwell's equations.  In Lorentz gauge
($\partial_\mu A^\mu = 0$) we must solve
\begin{equation}
  \partial^2 A^\mu = -j^\mu,
\end{equation}
or in Fourier space,
\begin{equation}
  \tilde{A}^\mu(k) = -\frac{1}{k^2}\tilde{j}^\mu(k) .
\end{equation}
Plugging in \eqref{eq:6.4}, we obtain a formula for the vector potential:
\begin{equation}
  A^\mu(x) = \int \frac{d^4k}{(2\pi)^4}\ e^{-ik\cdot x}\
    \frac{-ie}{k^2}\ \Bigl( \frac{p'^{\mu}}{k\cdot p' + i\epsilon}
      - \frac{p^{\mu}}{k\cdot p - i\epsilon} \Bigr).
  \label{eq:6.5}
\end{equation}
The $k^0$ integral can be performed as a contour integral in the complex
plane.  The locations of the poles are as follows: There are poles at
$k^0 = \pm|\mathbf{k}|$ (from the $1/k^2$ factor) and at
$k\cdot p = 0$ and $k\cdot p' = 0$.  We place the poles at
$k^0 = \pm|\mathbf{k}|$ below the real axis so that (as we shall soon confirm)
the radiation field will satisfy retarded boundary conditions.

For $t < 0$ we close the contour upward, picking up the pole at
$k\cdot p = 0$, that is, $k^0 = \mathbf{k}\cdot\mathbf{p}/p^0$.  The result is
\begin{equation}
  A^\mu(\mathbf{x}, t) \sim \int \frac{d^3k}{(2\pi)^3}\ e^{i\mathbf{k}\cdot\mathbf{x}}
    e^{-i(\mathbf{k}\cdot\mathbf{p}/p^0)t}\ (2\pi i)\Bigl( \frac{1}{2p^0}
      \cdot \frac{p^\mu}{-k^2}\cdots \Bigr).
\end{equation}
In the reference frame where the particle is initially at rest, its momentum
vector is $p^\mu = (p^0, 0)^\mu$, and the vector potential reduces to
\begin{equation}
  A^0(\mathbf{x}) = \int \frac{d^3k}{(2\pi)^3}\ e^{i\mathbf{k}\cdot\mathbf{x}}\
    \frac{e}{\mathbf{k}^2} = \frac{e}{4\pi|\mathbf{x}|} .
\end{equation}
This is just the Coulomb potential of an unaccelerated charge.  As we would
expect, there is no radiation field before the particle is scattered.

After scattering ($t > 0$), we close the contour downward, picking up the
three poles below the real axis.  The pole at
$k^0 = \mathbf{k}\cdot\mathbf{p}'/p'^0$ gives the Coulomb potential of the
outgoing particle.  Thus the other two poles are completely responsible for
the radiation field.  Their contribution gives
\begin{equation}
  A^\mu_{\mathrm{rad}}(\mathbf{x}, t) =
    \int \frac{d^3k}{(2\pi)^3}\ e^{ik\cdot x}\ \mathcal{A}^\mu(k),
  \label{eq:6.6}
\end{equation}
where the momentum-space amplitude $\mathcal{A}^\mu(k)$ is given by
\begin{equation}
  \mathcal{A}^\mu(k) = -ie\ \Bigl( \frac{p'^{\mu}}{k\cdot p'}
    - \frac{p^{\mu}}{k\cdot p} \Bigr).
  \label{eq:6.7}
\end{equation}
(The condition $k^0 = |\mathbf{k}|$ is implicit here and in the rest of this
calculation.)

To calculate the energy radiated, we must find the electric and magnetic
fields.  It is easiest to write $\mathbf{E}$ and $\mathbf{B}$ as the real parts
of complex Fourier integrals, just as we did for $A^\mu$:
\begin{equation}
  \mathbf{E}(\mathbf{x},t) = \mathrm{Re}\int \frac{d^3k}{(2\pi)^3}\,
    \bm{\mathcal{E}}(k)\,e^{ik\cdot x},\qquad
  \mathbf{B}(\mathbf{x},t) = \mathrm{Re}\int \frac{d^3k}{(2\pi)^3}\,
    \bm{\mathcal{B}}(k)\,e^{ik\cdot x},
\end{equation}
with
\begin{equation}
  \bm{\mathcal{E}}(k) = ik^0\bm{\mathcal{A}}(k)
    - i\mathbf{k}\mathcal{A}^0(k),\qquad
  \bm{\mathcal{B}}(k) = i\mathbf{k}\times\bm{\mathcal{A}}(k) .
\end{equation}
The energy radiated is
\begin{equation}
  \text{Energy} = \int dt \int d^2x\ \hat{\mathbf{n}}\cdot
    \bigl( \mathbf{E}\times\mathbf{B} \bigr) .
\end{equation}
Carrying out the integrals (see the details in the text), one finds
\begin{equation}
  \text{Energy} = \int \frac{d^3k}{(2\pi)^3}\ |\mathbf{k}|\,
    \Bigl( \frac{e^2}{4\pi} \Bigr)\ \mathcal{I}(\mathbf{v}, \mathbf{v}')
  = \frac{\alpha}{4\pi^2}\int_0^{\omega_{\mathrm{max}}} d\omega\,
    \omega \int d\Omega\ \mathcal{I}(\mathbf{v}, \mathbf{v}'),
  \label{eq:6.12}
\end{equation}
where $\mathcal{I}(\mathbf{v},\mathbf{v}')$ is a known function of the initial
and final electron velocities.

Now let us compute the function $\mathcal{I}(\mathbf{v},\mathbf{v}')$.
For nonrelativistic motion, $|\mathbf{v}| \ll 1$, and we can approximate
\begin{equation}
  \mathcal{I}(\mathbf{v},\mathbf{v}') \approx
    \frac{1}{4\pi}\ \Bigl( \frac{\mathbf{v}'}{1 - \hat{\mathbf{k}}\cdot\mathbf{v}'}
      - \frac{\mathbf{v}}{1 - \hat{\mathbf{k}}\cdot\mathbf{v}} \Bigr)^2 .
  \label{eq:6.13}
\end{equation}
The classical result for the radiation intensity is obtained by integrating
this expression.  The integral over angles is
\begin{equation}
  \int \frac{d\Omega}{4\pi}\ \Bigl( \frac{\mathbf{v}}{1-\hat{\mathbf{k}}\cdot\mathbf{v}}
    \Bigr)^2 = \frac{v^2}{\gamma^2(1-v^2)}
    = v^2\,\frac{1-v^2}{(1-v^2)^2}\ \cdots ,
\end{equation}
and the cross terms can be evaluated using
\begin{equation}
  \int \frac{d\Omega}{4\pi}\ \frac{1}{(1-\hat{\mathbf{k}}\cdot\mathbf{v})
    (1-\hat{\mathbf{k}}\cdot\mathbf{v}')}
  = \frac{1}{|\mathbf{v}-\mathbf{v}'|}\ \mathrm{arctanh}\,
    \frac{|\mathbf{v}-\mathbf{v}'|}{1 - \mathbf{v}\cdot\mathbf{v}'}
    \cdot\ \cdots .
  \label{eq:6.16}
\end{equation}
The integrals are easy to perform, and we obtain
\begin{equation}
  \mathcal{I}(\mathbf{v},\mathbf{v}') \approx
    \frac{1}{2}\Bigl( \frac{\mathbf{v}'^2}{1-v'^2} + \frac{\mathbf{v}^2}{1-v^2}
      \Bigr) - \frac{1}{2}\ \cdots ,
  \label{eq:6.17}
\end{equation}
where $q^2 = (p'-p)^2$.

In conclusion, we have found that the radiated energy at low frequencies is
given by
\begin{equation}
  \text{Energy} = \frac{\alpha}{\pi}\int_0^{\omega_{\mathrm{max}}}
    d\omega\ \mathcal{I}(\mathbf{v},\mathbf{v}')
  \xrightarrow{q^2 \gg m^2}
  \frac{\alpha}{\pi}\ \int_0^{\omega_{\mathrm{max}}} d\omega\
    \log\Bigl( \frac{q^2}{m^2} \Bigr).
  \label{eq:6.18}
\end{equation}
If this energy is made up of photons, each photon contributes energy $k$.  We
would then expect
\begin{equation}
  \text{Number of photons} = \frac{\alpha}{\pi}\ \int_0^{\omega_{\mathrm{max}}}
    \frac{dk}{k}\ \mathcal{I}(\mathbf{v},\mathbf{v}') .
  \label{eq:6.19}
\end{equation}
We hope that a quantum-mechanical calculation will confirm this result.

\subsection{Quantum Computation}

Consider now the quantum-mechanical process in which one photon is radiated
during the scattering of an electron:
\begin{equation}
  e^{-}(p) \to e^{-}(p') + \gamma(k) .
\end{equation}
Let $\mathcal{M}_0$ denote the part of the amplitude that comes from the
electron's interaction with the external field.  Then the amplitude for the
whole process is
\begin{equation}
  i\mathcal{M} = -ie\,\bar{u}(p')\ \Bigl[
    \gamma^\mu\epsilon^{*}_\mu(k)\ \frac{i(\slashed{p}'+\slashed{k}+m)}
      {(p'+k)^2-m^2}\ \mathcal{M}_0(p'+k, p)
    + \mathcal{M}_0(p', p-k)\ \frac{i(\slashed{p}-\slashed{k}+m)}
      {(p-k)^2-m^2}\ \gamma^\mu\epsilon^{*}_\mu(k) \Bigr] u(p) .
  \label{eq:6.20}
\end{equation}
Since we are interested in connecting with the classical limit, assume that
the photon radiated is soft: $|\mathbf{k}| \ll |\mathbf{p}'-\mathbf{p}|$.
Then we can approximate
\begin{equation}
  \mathcal{M}_0(p', p-k) \approx \mathcal{M}_0(p'+k, p)
    \approx \mathcal{M}_0(p', p),
  \label{eq:6.21}
\end{equation}
and we can ignore $\slashed{k}$ in the numerators of the propagators.  The
numerators can be further simplified with some Dirac algebra.  In the first
term we have
\begin{equation}
  (\slashed{p}'+m)\slashed{\epsilon}^{*}u(p) = \bigl[
    2p'\cdot\epsilon^{*} + \slashed{\epsilon}^{*}(\slashed{p}'-m) \bigr]u(p)
    = 2p'\cdot\epsilon^{*}\, u(p) .
\end{equation}
Similarly, in the second term,
\begin{equation}
  \bar{u}(p')\slashed{\epsilon}^{*}(\slashed{p}+m) = \bar{u}(p')\,2p\cdot\epsilon^{*}.
\end{equation}
The denominators of the propagators also simplify:
\begin{equation}
  (p-k)^2 - m^2 = -2p\cdot k;\qquad
  (p'+k)^2 - m^2 = 2p'\cdot k .
\end{equation}
So in the soft-photon approximation, the amplitude becomes
\begin{equation}
  i\mathcal{M} = \bar{u}(p')\mathcal{M}_0(p',p)u(p)\,
    (-ie)\ \Bigl( \frac{p'\cdot\epsilon^{*}}{p'\cdot k}
      - \frac{p\cdot\epsilon^{*}}{p\cdot k} \Bigr).
  \label{eq:6.22}
\end{equation}
This is just the amplitude for elastic scattering (without bremsstrahlung),
times a factor (in brackets) for the emission of the photon.

The cross section for our process is also easy to express in terms of the
elastic cross section; just insert an additional phase-space integration for
the photon variable $k$.  Summing over the two photon polarization states, we
have
\begin{equation}
  \begin{aligned}
    d\sigma(p\to p' + \gamma) &= d\sigma(p\to p')\cdot
      \int \frac{d^3k}{(2\pi)^3}\ \frac{1}{2k}\ \frac{\alpha}{2\pi}\\
    &\quad \times \sum_{\lambda=1,2}\ \Bigl|
      \frac{p'\cdot\epsilon^{(\lambda)}}{p'\cdot k}
      - \frac{p\cdot\epsilon^{(\lambda)}}{p\cdot k} \Bigr|^2 .
  \end{aligned}
  \label{eq:6.23}
\end{equation}
Thus the differential probability of radiating a photon with momentum $k$,
given that the electron scatters from $p$ to $p'$, is
\begin{equation}
  d(\mathrm{prob}) = \frac{\alpha}{(2\pi)^2}\ \frac{d^3k}{k}\
    \mathcal{I}(\mathbf{v},\mathbf{v}') .
  \label{eq:6.24}
\end{equation}
This looks very familiar; if we multiply by the photon energy $k$ to compute
the expected energy radiated, we recover the classical expression
\eqref{eq:6.12}.

But there is a problem.  Equation \eqref{eq:6.24} is an expression not for
the expected number of photons radiated, but for the probability of radiating
a single photon.  The problem becomes worse if we integrate over the photon
momentum.  As in \eqref{eq:6.16}, we can integrate only up to the energy at
which our soft-photon approximations break down; a reasonable estimate for
this energy is $|\mathbf{q}| = |\mathbf{p}-\mathbf{p}'|$.  The integral is
therefore
\begin{equation}
  \text{Total probability} \approx \frac{\alpha}{\pi}\int \frac{dk}{k}\
    \mathcal{I}(\mathbf{v}, \mathbf{v}') .
  \label{eq:6.25}
\end{equation}
Since $\mathcal{I}(\mathbf{v},\mathbf{v}')$ is independent of $k$, the
integral diverges at its lower limit (where all our approximations are well
justified).  In other words, the total probability of radiating a very soft
photon is infinite.  This is the famous problem of \emph{infrared divergences}
in QED perturbation theory.

We can artificially make the integral in \eqref{eq:6.25} well-defined by
pretending that the photon has a very small mass $\mu$.  This mass would then
provide a lower cutoff for the integral, allowing us to write the result of
this section as
\begin{equation}
  \begin{aligned}
    d\sigma(p\to p'+\gamma(k)) &= d\sigma(p\to p')\cdot
      \frac{\alpha}{2\pi}\log\Bigl( \frac{-q^2}{\mu^2} \Bigr)
      \mathcal{I}(\mathbf{v},\mathbf{v}')\\
    &\xrightarrow{-q^2\gg m^2}
    d\sigma(p\to p')\cdot \frac{\alpha}{2\pi}
      \log\Bigl( \frac{-q^2}{\mu^2} \Bigr)\log\Bigl( \frac{-q^2}{m^2} \Bigr).
  \end{aligned}
  \label{eq:6.26}
\end{equation}
The $q^2$ dependence of this result, known as the \emph{Sudakov double
logarithm}, is physical and will appear again in Section~6.4.  The dependence
on $\mu$, however, presents a problem that we must solve.  It is not hard to
guess that the resolution of this problem will involve reinterpreting
\eqref{eq:6.24} as the expected number of radiated photons, rather than the
probability of radiating a single photon.  We will see in Sections~6.4 and~6.5
how this reinterpretation follows from the Feynman diagrams.  To prepare for
that discussion, however, we need to improve our understanding of the
amplitude for scattering without radiation.

\section{The Electron Vertex Function: Formal Structure}

Having briefly discussed QED radiative corrections due to emission of photons
(bremsstrahlung), let us now study the correction to electron scattering that
comes from the presence of an additional virtual photon:
\begin{equation}
  \text{(one-loop vertex diagram)} .
  \label{eq:6.27}
\end{equation}
This will be our first experience with a Feynman diagram containing a loop.
Such diagrams give rise to significant and profound complications in quantum
field theory.

The result of computing this diagram will be rather complicated, so it will be
useful to think ahead about what form we expect this correction to take and
how to interpret its various possible terms.  In this section, we will
consider the general properties of vertex correction diagrams.  We will see
that the basic requirements of Lorentz invariance, the discrete symmetries of
QED, and the Ward identity strongly constrain the form of the vertex.

Consider, then, the class of diagrams
\begin{equation}
  \text{(electron-photon vertex)} = -ie\Gamma^\mu(p', p),
\end{equation}
where the gray circle indicates the sum of the lowest-order electron-photon
vertex and all amputated loop corrections.  We will call this sum of vertex
diagrams $-ie\Gamma^\mu(p',p)$.  Then, according to our master formula
\eqref{eq:4.103} for $S$-matrix elements, the amplitude for electron scattering
from a heavy target is
\begin{equation}
  i\mathcal{M} = ie^2\ \bigl( \bar{u}(p')\Gamma^\mu(p',p)u(p) \bigr)\
    \frac{1}{q^2}\ \bigl( \bar{u}(k')\gamma_\mu u(k) \bigr).
  \label{eq:6.28}
\end{equation}
More generally, the function $\Gamma^\mu(p',p)$ appears in the $S$-matrix
element for the scattering of an electron from an external electromagnetic
field.  As in Problem~4.4, add to the Hamiltonian of QED the interaction
\begin{equation}
  \Delta H_{\mathrm{int}} = \int d^3x\ e\,j^\mu(x)\,A^{cl}_\mu(x),
  \label{eq:6.29}
\end{equation}
where $j^\mu(x) = \bar{\psi}(x)\gamma^\mu\psi(x)$ is the electromagnetic
current and $A^{cl}_\mu$ is a fixed classical potential.  In the leading order
of perturbation theory, the $S$-matrix element for scattering from this field
is
\begin{equation}
  i\mathcal{M}(2\pi)\delta(p^{0'}-p^0) = -ie\,\bar{u}(p')\gamma^\mu u(p)\,
    \tilde{A}^{cl}_\mu(p'-p),
\end{equation}
where $\tilde{A}^{cl}_\mu(q)$ is the Fourier transform of $A^{cl}_\mu(x)$.
The vertex corrections modify this expression to
\begin{equation}
  i\mathcal{M}(2\pi)\delta(p^{0'}-p^0) = -ie\,\bar{u}(p')\Gamma^\mu(p',p)u(p)\,
    \tilde{A}^{cl}_\mu(p'-p) .
  \label{eq:6.30}
\end{equation}
In writing \eqref{eq:6.28} and \eqref{eq:6.30}, we have deliberately omitted
the contribution of vacuum polarization diagrams, such as the fourth diagram
of \eqref{eq:6.1}.  The reason for this omission is that these diagrams should
be considered corrections to the photon propagator, not to the electron-photon
vertex.

Since $\Gamma^\mu(p',p)$ is a $4\times4$ matrix in Dirac space, we can express
it as a linear combination of the sixteen basis matrices of Section~3.4.  The
requirements of Lorentz invariance and the discrete symmetries $C$, $P$, and
$T$ restrict the allowed terms.  One finds that the most general form
consistent with these symmetries is
\begin{equation}
  \Gamma^\mu(p',p) = \gamma^\mu\,F_1(q^2)
    + \frac{i\sigma^{\mu\nu}q_\nu}{2m}\,F_2(q^2),
  \label{eq:6.33}
\end{equation}
where $q = p'-p$, $F_1$ and $F_2$ are scalar functions of $q^2$ (with
$F_1(0) = 1$), and the $F_2$ term is the \emph{anomalous magnetic moment}
piece.  To see this, use the Gordon identity
\eqref{eq:6.32},
\begin{equation}
  \bar{u}(p')\gamma^\mu u(p) = \bar{u}(p')\Bigl(
    \frac{p'^{\mu}+p^{\mu}}{2m} + \frac{i\sigma^{\mu\nu}q_\nu}{2m} \Bigr)u(p),
  \label{eq:6.32}
\end{equation}
which shows that the $\sigma^{\mu\nu}$ term in \eqref{eq:6.33} contributes an
additional magnetic moment.  The constraint that $\Gamma^\mu$ must satisfy the
Ward identity
\begin{equation}
  q_\mu\Gamma^\mu(p',p) = 0
\end{equation}
implies that $F_2$ is completely unconstrained by the Ward identity, while
$F_1$ satisfies $F_1(0) = 1$ (after renormalization).

\section{The Electron Vertex Function: Evaluation}

Now that we know what form the answer is to take (Eq.~\eqref{eq:6.33}), we are
ready to evaluate the one-loop contribution to the electron vertex function.
Assign momenta on the diagram as follows: the incoming electron has momentum
$p$, the outgoing electron has momentum $p'$, and the virtual photon has
momentum $k-p$ (so that $q = p'-p$ is the momentum transfer).  Applying the
Feynman rules, we find, to order $\alpha$, that
$\Gamma^\mu = \gamma^\mu + \delta\Gamma^\mu$, where
\begin{equation}
  \begin{aligned}
    \bar{u}(p')\,\delta\Gamma^\mu(p',p)\,u(p)
    &= \int \frac{d^4k}{(2\pi)^4}\
      \frac{-ig_{\nu\rho}}{(k-p)^2+i\epsilon}\ \
      \bar{u}(p')(-ie\gamma^\nu)\
      \frac{i(\slashed{k}'+m)}{k'^2-m^2+i\epsilon}\
      \frac{i(\slashed{k}+m)}{k^2-m^2+i\epsilon}\
      (-ie\gamma^\rho)\,u(p)\\
    &= 2e^2\int \frac{d^4k}{(2\pi)^4}\
      \frac{\bar{u}(p')\bigl[ \gamma^\mu k^2 + m^2\gamma^\mu
        - 2m(k+k')^\mu \bigr]u(p)}
      {((k-p)^2+i\epsilon)(k'^2-m^2+i\epsilon)(k^2-m^2+i\epsilon)} .
  \end{aligned}
  \label{eq:6.38}
\end{equation}
In the second line we have used the contraction identity
$\gamma^\nu\gamma^\mu\gamma_\nu = -2\gamma^\mu$.  Note that the $+i\epsilon$
terms in the denominators cannot be dropped; they are necessary for proper
evaluation of the loop-momentum integral.

The integral looks impossible, and in fact it will not be easy.  The
evaluation of such integrals requires another piece of computational
technology, known as the \emph{method of Feynman parameters} (although a very
similar method was introduced earlier by Schwinger).

\subsection{Feynman Parameters}

The goal of this method is to squeeze the three denominator factors of
\eqref{eq:6.38} into a single quadratic polynomial in $k$, raised to the third
power.  We can then shift $k$ by a constant to complete the square in this
polynomial and evaluate the remaining spherically symmetric integral without
difficulty.  The price will be the introduction of auxiliary parameters to be
integrated over.

It is easiest to begin with the simpler case of two factors in the
denominator.  We would then use the identity
\begin{equation}
  \frac{1}{AB} = \int_0^1 dx\ \frac{1}{[xA+(1-x)B]^2}
    = \int dx\,dy\ \delta(x+y-1)\ \frac{1}{[xA+yB]^2} .
  \label{eq:6.39}
\end{equation}
An example of its use might look like this:
\begin{equation}
  \frac{1}{(k-p)^2(k^2-m^2)}
  = \int dx\,dy\ \delta(x+y-1)\
    \frac{1}{[k^2 - 2xk\cdot p + xp^2 - ym^2]^2} .
\end{equation}
If we now let $\ell = k - xp$, we see that the denominator depends only on
$\ell^2$.  Integrating over $d^4k$ would now be much easier, since
$d^4k = d^4\ell$ and the integrand is spherically symmetric with respect to
$\ell$.  The variables $x$ and $y$ that make this transformation possible are
called \emph{Feynman parameters}.

Our integral \eqref{eq:6.38} involves a denominator with three factors, so we
need a slightly better identity.  By differentiating \eqref{eq:6.39} with
respect to $B$, it is easy to prove
\begin{equation}
  \frac{1}{AB^n} = \int dx\,dy\ \delta(x+y-1)\
    \frac{n\,y^{n-1}}{[xA+yB]^{n+1}} .
  \label{eq:6.40}
\end{equation}
But this still isn't quite good enough.  The formula we need is
\begin{equation}
  \frac{1}{A_1A_2\cdots A_n}
  = \int_0^1 dx_1\cdots dx_n\ \delta\Bigl( \textstyle\sum x_i - 1 \Bigr)\
    \frac{(n-1)!}{[x_1A_1 + x_2A_2 + \cdots + x_nA_n]^n} .
  \label{eq:6.41}
\end{equation}
The proof of this identity is by induction.  The case $n=2$ is just
Eq.~\eqref{eq:6.39}; the induction step is not difficult and involves the use
of \eqref{eq:6.40}.

By repeated differentiation of \eqref{eq:6.41}, you can derive the even more
general identity
\begin{equation}
  \frac{1}{A_1^{m_1}A_2^{m_2}\cdots A_n^{m_n}}
  = \int_0^1 dx_1\cdots dx_n\ \delta\Bigl( \textstyle\sum x_i - 1 \Bigr)\
    \frac{\Gamma(m_1+\cdots+m_n)}{\Gamma(m_1)\cdots\Gamma(m_n)}\
    \frac{x_1^{m_1-1}\cdots x_n^{m_n-1}}
      {[x_1A_1+\cdots+x_nA_n]^{m_1+\cdots+m_n}} .
  \label{eq:6.42}
\end{equation}
This formula is true even when the $m_i$ are not integers; in Section~10.5 we
will apply it in such a case.

\subsection{Evaluation of the Form Factors}

Now let us apply formula \eqref{eq:6.41} to the denominator of
\eqref{eq:6.38}:
\begin{equation}
  \frac{1}{((k-p)^2+i\epsilon)(k'^2-m^2+i\epsilon)(k^2-m^2+i\epsilon)}
  = \int dx\,dy\,dz\ \delta(x+y+z-1)\ \frac{2}{D^3} ,
\end{equation}
where the new denominator $D$ is
\begin{equation}
  \begin{aligned}
    D &= x(k^2-m^2) + y(k'^2-m^2) + z(k-p)^2 + (x+y+z)i\epsilon\\
    &= k^2 + 2k\cdot(yq-zp) + yq^2 + zp^2 - (x+y)m^2 + i\epsilon .
  \end{aligned}
\end{equation}
In the second line we have used $x+y+z=1$ and $k' = k+q$.  Now shift $k$ to
complete the square:
\begin{equation}
  \ell = k + yq - zp .
\end{equation}
After a bit of algebra we find that $D$ simplifies to
\begin{equation}
  D = \ell^2 - \Delta + i\epsilon,
  \label{eq:6.43}
\end{equation}
where
\begin{equation}
  \Delta = -xyq^2 + (1-z)^2m^2 .
  \label{eq:6.44}
\end{equation}
Since $q^2 < 0$ for a scattering process, $\Delta$ is positive; we can think
of it as an effective mass term.

Next we must express the numerator of \eqref{eq:6.38} in terms of $\ell$.
This task is simplified by noting that since $D$ depends only on the
magnitude of $\ell$,
\begin{equation}
  \int \frac{d^4\ell}{(2\pi)^4}\ \frac{\ell^\mu}{D^3} = 0;
  \label{eq:6.45}
\end{equation}
\begin{equation}
  \int \frac{d^4\ell}{(2\pi)^4}\ \frac{\ell^\mu\ell^\nu}{D^3}
    = \frac{g^{\mu\nu}}{4}\int \frac{d^4\ell}{(2\pi)^4}\ \frac{\ell^2}{D^3} .
  \label{eq:6.46}
\end{equation}
The first identity follows from symmetry.  To prove the second, note that the
integral vanishes by symmetry unless $\mu = \nu$.  Lorentz invariance therefore
requires that we get something proportional to $g^{\mu\nu}$.  To check the
coefficient, contract each side with $g_{\mu\nu}$.  Using these identities, we
have
\begin{equation}
  \begin{aligned}
    \text{Numerator} &= \bar{u}(p')\Bigl[ \gamma^\mu\ell^2 + m^2\gamma^\mu
      - 2m(k+k')^\mu \Bigr]u(p)\\
    &\to \bar{u}(p')\Bigl[ -\tfrac{1}{2}\gamma^\mu\ell^2
      + \cdots \Bigr]u(p),
  \end{aligned}
\end{equation}
(Remember that $k' = k+q$.)
Putting the numerator into a useful form is now just a matter of some tedious
Dirac algebra (about a page or two).  This is where our work in the last
section pays off, since it tells us what kind of an answer to expect.  We
eventually want to group everything into two terms, proportional to
$\gamma^\mu$ and $i\sigma^{\mu\nu}q_\nu$.  The most straightforward way to
accomplish this is to aim instead for an expression of the form
\begin{equation}
  \gamma^\mu A + (p'^{\mu}+p^\mu)B + q^\mu C,
\end{equation}
just as in Eq.~(6.31).  Attaining this form requires only the
anticommutation relations (for example, $\slashed{p}\gamma^\mu =
2p^\mu - \gamma^\mu\slashed{p}$) and the Dirac equation $\slashed{p}u(p) =
mu(p)$ and $\bar{u}(p')\slashed{p}' = \bar{u}(p')m$; note that this implies
$\bar{u}(p')\slashed{q}u(p) = 0$.  It is also useful to remember that
$x+y+z = 1$.  When the smoke clears, we have
\begin{equation}
  \begin{aligned}
    \text{Numerator} &= \bar{u}(p')\Bigl[ \gamma^\mu\,
      \bigl( \tfrac{1}{2}\ell^2 + (1-x)(1-y)q^2 + (1-2z-z^2)m^2 \bigr)\\
    &\quad + (p'^{\mu}+p^\mu)\cdot mz(z-1)
      + q^\mu\cdot m(z-2)(x-y) \Bigr]u(p) .
  \end{aligned}
\end{equation}
The coefficient of $q^\mu$ must vanish according to the Ward identity, as
discussed after Eq.~(6.31).  To see that it does, note from
\eqref{eq:6.44} that the denominator is symmetric under $x\leftrightarrow y$.
The coefficient of $q^\mu$ is odd under $x\leftrightarrow y$ and therefore
vanishes when integrated over $x$ and $y$.

Still following our work in the previous section, we now use the Gordon
identity \eqref{eq:6.32} to eliminate $(p'+p)$ in favor of
$i\sigma^{\mu\nu}q_\nu$.  Our entire expression for the $O(\alpha)$
contribution to the electron vertex then becomes
\begin{equation}
  \begin{aligned}
    \bar{u}(p')\,\delta\Gamma^\mu(p',p)\,u(p)
    &= 2ie^2 \int dx\,dy\,dz\ \delta(x+y+z-1)\ \int
      \frac{d^4\ell}{(2\pi)^4}\ \frac{1}{D^3}\\
    &\quad \times \bar{u}(p')\Bigl[
      \gamma^\mu\Bigl( \tfrac{1}{2}\ell^2 + (1-x)(1-y)q^2
        + (1-4z+z^2)m^2 \Bigr)\\
      &\quad + \frac{i\sigma^{\mu\nu}q_\nu}{2m}\,
        \bigl( 2m^2z(1-z) \bigr) \Bigr]u(p),
  \end{aligned}
  \label{eq:6.47}
\end{equation}
where as before,
\begin{equation}
  D = \ell^2 - \Delta + i\epsilon,\qquad
  \Delta = -xyq^2 + (1-z)^2m^2 > 0 .
\end{equation}
The decomposition into form factors is now manifest.

With most of the work behind us, our main remaining task is to perform the
momentum integral.  It is not difficult to evaluate the $\ell^0$ integral as a
contour integral, then perform the remaining three-dimensional integral by
Wick rotation.  The rotation of the $\ell^0$ contour is shown in
Fig.~\ref{fig:6.1}.

\begin{figure}[t]
  \centering
  \includegraphics[width=0.55\textwidth]{fig6.1.png}
  \caption{The contour of the $\ell^0$ integration can be rotated as shown.}
  \label{fig:6.1}
\end{figure}

The result of the momentum integral is
\begin{equation}
  \int \frac{d^4\ell}{(2\pi)^4}\ \frac{1}{(\ell^2-\Delta+i\epsilon)^3}
    = \frac{-i}{32\pi^2\Delta},
  \qquad
  \int \frac{d^4\ell}{(2\pi)^4}\ \frac{\ell^2}{(\ell^2-\Delta+i\epsilon)^3}
    = \frac{i}{16\pi^2}\ \cdots ,
  \label{eq:6.48}
\end{equation}
which, when inserted into \eqref{eq:6.47}, yields the one-loop form factors.
In the ultraviolet, the integrals diverge logarithmically; we regulate them
with a cutoff $\Lambda$ (or dimensional regularization).  With both of these
provisional modifications, the form factors are
\begin{equation}
  F_1(q^2) = 1 + \frac{\alpha}{2\pi}\int dx\,dy\,dz\
    \delta(x+y+z-1)\ \log\Bigl(
      \frac{m^2(1-z)^2}{m^2(1-z)^2 - q^2xy} \Bigr)\
    + \frac{\alpha}{2\pi}\int dx\,dy\,dz\ \delta(x+y+z-1)\
      \frac{m^2(1-4z+z^2) + q^2(1-x)(1-y)}
        {m^2(1-z)^2 - q^2xy + \mu^2 z}\
    + \mathcal{O}(\alpha^2);
  \label{eq:6.56}
\end{equation}
\begin{equation}
  F_2(q^2) = \frac{\alpha}{2\pi}\int dx\,dy\,dz\
    \delta(x+y+z-1)\ \frac{2m^2z(1-z)}{m^2(1-z)^2 - q^2xy}
    + \mathcal{O}(\alpha^2) .
  \label{eq:6.57}
\end{equation}
Note that neither the ultraviolet nor the infrared divergence affects
$F_2(q^2)$.  We can therefore evaluate unambiguously
\begin{equation}
  F_2(q^2 = 0) = \frac{\alpha}{2\pi}\int_0^1 dz\ \int_0^{1-z} dy\
    \frac{2m^2z(1-z)}{m^2(1-z)^2}
    = \frac{\alpha}{2\pi}\ \int_0^1 dz\ \frac{z}{1-z}\ \cdots
    = \frac{\alpha}{2\pi} .
  \label{eq:6.58}
\end{equation}
Thus, we get a correction to the $g$-factor of the electron:
\begin{equation}
  \frac{g-2}{2} = \frac{\alpha}{2\pi} = 0.0011614 .
  \label{eq:6.59}
\end{equation}
This result was first obtained by Schwinger in 1948.%
\footnote{J.~Schwinger, \emph{Phys.\ Rev.} \textbf{73}, 416L (1948).}
Experiments give $a_e = 0.0011597$.  Apparently, the unambiguous value that we
obtained for $F_2(0)$ is also, up to higher orders in $\alpha$, unambiguously
correct.

\subsection{Precision Tests of QED}

Building on the success of the order-$\alpha$ QED prediction for $a_e$,
successive generations of physicists have improved the accuracy of both the
theoretical and the experimental determination of this quantity.  The
coefficients of the QED formula for $a_e$ are now known through order
$\alpha^4$.  The calculation of the order-$\alpha^2$ and higher coefficients
requires a systematic treatment of ultraviolet divergences.

These challenging theoretical calculations have been matched by increasingly
imaginative experiments.  The most recent measurement of $a_e$ uses a
technique, developed by Dehmelt and collaborators, in which individual
electrons are trapped in a system of electrostatic and magnetostatic fields
and excited to a spin resonance.%
\footnote{R.~Van Dyck, Jr., P.~Schwinberg, and H.~Dehmelt, \emph{Phys.\ Rev.\
Lett.} \textbf{59}, 26 (1987).}  Today, the best theoretical and experimental
values of $a_e$ agree to eight significant figures.

High-order QED calculations have also been carried out for several other
quantities.  These include transition energies in hydrogen and hydrogen-like
atoms, the anomalous magnetic moment of the muon, and the decay rates of
singlet and triplet positronium.  Many of these quantities have also been
measured to high precision.  The full set of these comparisons gives a detailed
test of the validity of QED in a variety of settings.  The results of these
precision tests are summarized in Table~\ref{tab:6.1}.

There is some subtlety in reporting the results of precision comparisons
between QED theory and experiment, since theoretical predictions require an
extremely precise value of $\alpha$, which can only be obtained from another
precision QED experiment.  We therefore quote each comparison between theory
and experiment as an independent determination of $\alpha$.  Each value of
$\alpha$ is assigned an error that is the composite of the expected
uncertainties from theory and experiment.  QED is confirmed to the extent that
the values of $\alpha$ from different sources agree.

The first nine entries in Table~\ref{tab:6.1} refer to QED calculations in
atomic physics settings.  Of these, the hydrogen hyperfine splitting, measured
using Ramsey's hydrogen maser, is the most precisely known quantity in
physics.  Unfortunately, the influence of the internal structure of the proton
leads to uncertainties that limit the accuracy with which this quantity can
be predicted theoretically.  The same difficulty applies to the Lamb shift,
the splitting between the $j = 1/2$ $2S$ and $2P$ levels of hydrogen.  The
most accurate QED tests now come from systems that involve no strongly
interacting particles, the electron $g-2$ and the hyperfine splitting in the
$e^{-}\mu^{+}$ atom, muonium.  The last entry in this group gives a new method
for determining $\alpha$, by converting a very accurate measurement of the
neutron Compton wavelength, using accurately known mass ratios, to a value of
the electron mass.  This can be combined with the known value of the Rydberg
energy and accurate QED formulae to determine $\alpha$.  The only serious
discrepancy among these numbers comes in the triplet positronium decay rate;
however, there is some evidence that diagrams of relative order $\alpha^2$
give a large correction to the value quoted in the table.

The next two entries are determinations of $\alpha$ from higher-order QED
reactions at high-energy electron colliders.  These high-energy experiments
typically achieve only percent-level accuracy, but their results are
consistent with the precise information available at lower energies.

Finally, the last two entries in the table give two independent measurements
of $\alpha$ from exotic quantum interference phenomena in condensed-matter
systems.  These two effects provide a standard resistance and a standard
frequency, respectively, which are believed to measure the charge of the
electron with corrections that are strictly zero for macroscopic systems.%
\footnote{For a discussion of these effects, and their exact relation to
$\alpha$, see D.~R.~Yennie, \emph{Rev.\ Mod.\ Phys.} \textbf{59}, 781 (1987).}

\begin{table}[t]
  \centering
  \caption{Values of $1/\alpha$ Obtained from Precision QED Experiments}
  \begin{tabular}{ll}
    \toprule
    Low-Energy QED: & \\
    Electron $(g-2)$ & 137.035\,992\,35\,(73)\\
    Muon $(g-2)$ & 137.035\,5\,(11)\\
    Muonium hyperfine splitting & 137.035\,994\,(18)\\
    Lamb shift & 137.036\,8\,(7)\\
    Hydrogen hyperfine splitting & 137.036\,0\,(3)\\
    $2\,^3S_1 - 1\,^3S_1$ splitting in positronium & 137.034\,(16)\\
    $1\,S_0$ positronium decay rate & 137.00\,(6)\\
    $3\,S_1$ positronium decay rate & 136.971\,(6)\\
    Neutron Compton wavelength & 137.036\,010\,1\,(54)\\
    \midrule
    High-Energy QED: & \\
    $\sigma(e^{+}e^{-}\to e^{+}e^{-}e^{+}e^{-})$ & 136.5\,(2.7)\\
    $\sigma(e^{+}e^{-}\to e^{+}e^{-}\mu^{+}\mu^{-})$ & 139.9\,(1.2)\\
    \midrule
    Condensed Matter: & \\
    Quantum Hall effect & 137.035\,997\,9\,(32)\\
    AC Josephson effect & 137.035\,977\,0\,(77)\\
    \bottomrule
  \end{tabular}
  \label{tab:6.1}
\end{table}

Each value of $\alpha$ displayed in this table is obtained by fitting an
experimental measurement to a theoretical expression that contains $\alpha$ as
a parameter.  The numbers in parentheses are the standard errors in the last
displayed digits, including both theoretical and experimental uncertainties.
This table is based on results presented in the survey of precision QED of
Kinoshita (1990).  That book contains a series of lucid reviews of the
remarkable theoretical and experimental technology that has been developed
for the detailed analysis of QED processes.  The five most accurate values are
updated as given by T.~Kinoshita in \emph{History of Original Ideas and Basic
Discoveries in Particle Physics}, H.~Newman and T.~Ypsilantis, eds.\ (Plenum
Press, New York, 1995).  This latter paper also gives an interesting
perspective on the future of precision QED experiments.

The entire picture fits together well beyond any reasonable expectation.  On
the evidence presented in this table, QED is the most stringently tested---and
the most dramatically successful---of all physical theories.

\section{The Electron Vertex Function: Infrared Divergence}

Now let us confront the infrared divergence in our result \eqref{eq:6.56} for
$F_1(q^2)$.  The dominant part, in the $\mu\to 0$ limit, is
\begin{equation}
  F_1(q^2) \sim \frac{\alpha}{2\pi}\int dx\,dy\,dz\ \delta(x+y+z-1)\
    \frac{m^2(1-4z+z^2) + q^2(1-x)(1-y)}
      {m^2(1-z)^2 - q^2xy + \mu^2z}
    + \cdots .
  \label{eq:6.60}
\end{equation}
To understand this expression we must do some work to simplify it, extracting
and evaluating the divergent part of the integral.  Throughout this section we
will retain only terms that diverge in the limit $\mu\to 0$.

First note that the divergence occurs in the corner of Feynman-parameter space
where $z\approx 1$ (and therefore $x\approx y\approx 0$).  In this region we
can set $z = 1$ and $x = y = 0$ in the numerators of \eqref{eq:6.60}.  We can
also set $z = 1$ in the $\mu^2$ terms in the denominators.  Using the delta
function to evaluate the $x$-integral, we then have
\begin{equation}
  F_1(q^2) = \frac{\alpha}{2\pi}\int_0^1 dz\ \int_0^{1-z} dy\
    \Bigl[ \frac{-2m^2 + q^2}{m^2(1-z)^2 - q^2y(1-z-y) + \mu^2}
      + \frac{m^2}{(1-z)^2 + \mu^2/m^2} \Bigr].
\end{equation}
(The lower limit on the $z$-integral is unimportant.)  Making the variable
changes
\begin{equation}
  v = 1 - z,\qquad w = \frac{y}{1-z},
\end{equation}
this expression becomes
\begin{equation}
  F_1(q^2) \sim \frac{\alpha}{2\pi}\int_0^1 dv\
    \int_0^1 dw\ \Bigl[ \frac{-2m^2}{m^2w^2 + \mu^2}
      + \frac{m^2}{m^2v^2 + \mu^2}\ \cdots \Bigr],
\end{equation}
\begin{equation}
  \sim \frac{\alpha}{2\pi}\ \Bigl[ \cdots \Bigr]
    \log\Bigl( \frac{-q^2}{\mu^2} \Bigr)\log\Bigl( \frac{-q^2}{m^2}
      \Bigr) + \cdots ,
  \label{eq:6.61}
\end{equation}
which is the Sudakov double logarithm found in Section~6.1.  The physical
requirement that this divergence cancel against the bremsstrahlung
contribution determines the subtraction and the physical interpretation of
the result.

\section{Summation and Interpretation of Infrared Divergences}

Infrared divergences arise from photons with ``soft'' momenta: real photons
with energy less than some cutoff $E_{-}$, and virtual photons with (after
Wick rotation) $k^2 < E_{-}^2$.  A typical higher-order diagram will involve
numerous real and virtual photons.  But to find a divergence, we need more
than a soft photon; we need a singular denominator in an electron propagator.
Consider, for example, the following two diagrams:
\begin{equation}
  \text{(soft photon emitted before hard photon)} \qquad
  \text{(soft photon emitted after hard photon)} .
\end{equation}
The first diagram, in which the electron emits a soft photon followed by a
hard photon, has no infrared divergence, since the momenta in both electron
propagators are far from the mass shell.  If the soft photon is emitted last,
however, the denominator of the adjacent propagator is
$(p'+k)^2 - m^2 = 2p'\cdot k$, which vanishes as $k\to 0$.  Thus the second
diagram does contain a divergence.  We would like, then, to consider diagrams
in which an arbitrary hard process, possibly involving emission of hard and
soft photons, is modified by the addition of soft real and virtual photons on
the electron legs:
\begin{equation}
  \underbrace{\text{soft real and virtual photons}}_{\text{on the electron
    legs}}\quad \underbrace{\text{arbitrary hard process}} .
\end{equation}
Following Weinberg, we will add up the contributions of all such diagrams.
The only new difficulty in this calculation will be in the combinatorics of
counting all the ways in which the photons can appear.

First consider the outgoing electron line.  We attach $n$ photons to the line,
with momenta $k_1, \ldots, k_n$.  For the moment we do not care whether these
are external photons, virtual photons connected to each other, or virtual
photons connected to vertices on the incoming electron line.  The Dirac
structure of this diagram is
\begin{equation}
  \bar{u}(p')\,(-ie\gamma^{\mu_1})\,
  \frac{i(\slashed{p}'+\slashed{k}_1+m)}{2p'\cdot k_1 + \mathcal{O}(k^2)}\,
  (-ie\gamma^{\mu_2})\,
  \frac{i(\slashed{p}'+\slashed{k}_1+\slashed{k}_2+m)}
    {2p'\cdot(k_1+k_2) + \mathcal{O}(k^2)}\ \cdots\,
  (-ie\gamma^{\mu_n})\,
  \frac{i(\slashed{p}'+\textstyle\sum\slashed{k}_i+m)}
    {2p'\cdot(k_1+\cdots+k_n) + \mathcal{O}(k^2)}\,
  \mathcal{M}_{\mathrm{hard}} .
  \label{eq:6.73}
\end{equation}
We will assume that all the $k_i$ are small, dropping the
$\mathcal{O}(k^2)$ terms in the denominators.  We will also drop the
$\slashed{k}$ terms in the numerators, just as in our treatment of
bremsstrahlung in Section~6.1.  Also, as we did there, we can push the factors
of $(\slashed{p}'+m)$ to the left and use
$\bar{u}(p')(\slashed{p}'-m) = 0$:
\begin{equation}
  \bar{u}(p')\slashed{\epsilon}^{\mu_1}(\slashed{p}'+m)
    \slashed{\epsilon}^{\mu_2}(\slashed{p}'+m)\cdots
  = \bar{u}(p')\,2p'\cdot\epsilon^{\mu_1}\slashed{\epsilon}^{\mu_2}
    (\slashed{p}'+m) = \bar{u}(p')\,2p'\cdot\epsilon^{\mu_1}\,
    2p'\cdot\epsilon^{\mu_2}\cdots .
\end{equation}
This turns expression \eqref{eq:6.73} into
\begin{equation}
  \bar{u}(p')\,(-ie)^n\ \prod_{i=1}^{n}\ \frac{p'\cdot\epsilon^{\mu_i}}
    {p'\cdot(k_1+\cdots+k_i)}\ \cdots .
  \label{eq:6.74}
\end{equation}
Still working with only the outgoing electron line, we must now sum over all
possible orderings of the momenta $k_1\cdots k_n$.  (This procedure will
overcount when two of the photons are attached together to form a single
virtual photon.  We will deal with this overcounting later.)  There are $n!$
different diagrams to sum, corresponding to the $n!$ permutations of the $n$
photon momenta.  Let $\pi$ denote one such permutation, so that $\pi(i)$ is
the number between 1 and $n$ that $i$ is taken to.  (For example, if $\pi$
denotes the permutation that takes $1\to 3$, $2\to 1$ and $3\to 2$, then
$\pi(1) = 3$, $\pi(2) = 1$, and $\pi(3) = 2$.)

Armed with this notation, we can perform the sum over permutations by means of
the following identity:
\begin{equation}
  \sum_{\text{all permutations }\pi}\ \frac{1}{p\cdot k_{\pi(1)}}\,
    \frac{1}{p\cdot(k_{\pi(1)}+k_{\pi(2)})}\ \cdots\,
    \frac{1}{p\cdot(k_{\pi(1)}+\cdots+k_{\pi(n)})}
  = \frac{1}{p\cdot k_1\,p\cdot k_2\cdots p\cdot k_n} .
  \label{eq:6.75}
\end{equation}
The proof of this formula proceeds by induction on $n$.  For $n=2$ we have
\begin{equation}
  \frac{1}{p\cdot k_{\pi(1)}}\frac{1}{p\cdot(k_{\pi(1)}+k_{\pi(2)})}
  = \frac{1}{p\cdot k_1\,p\cdot k_2},
\end{equation}
which is easily verified by adding the two terms.  For the induction step,
notice that the last factor on the left-hand side of \eqref{eq:6.75} is the
same for every permutation $\pi$.  Pulling this factor outside the sum, the
left-hand side becomes
\begin{equation}
  \text{LHS} = \sum_{\pi}\ \frac{1}{p\cdot k_{\pi(1)}}\ \cdots
    \frac{1}{p\cdot(k_{\pi(1)}+\cdots+k_{\pi(n-1)})}\ .
\end{equation}
For any given $\pi$, the quantity being summed is independent of
$k_{\pi(n)}$.  Letting $i = \pi(n)$, we can now write
\begin{equation}
  \sum_{\pi} = \sum_{i=1}^{n}\ \sum_{\pi'(i)},
\end{equation}
where $\pi'(i)$ is the set of all permutations on the remaining $n-1$
integers.  Assuming by induction that \eqref{eq:6.75} is true for $n-1$, we
have
\begin{equation}
  \text{LHS} = \frac{1}{p\cdot k_1\,p\cdot k_2\cdots p\cdot k_{i-1}\,
    p\cdot k_{i+1}\cdots p\cdot k_n}\ .
\end{equation}
If we now multiply and divide each term in this sum by $p\cdot k_i$, we easily
obtain our desired result \eqref{eq:6.75}.

Applying \eqref{eq:6.75} to \eqref{eq:6.74}, we find
\begin{equation}
  \bar{u}(p')\,(-ie)^n\ \Bigl( \frac{p'\cdot\epsilon_1}{p'\cdot k_1}
    \Bigr)\Bigl( \frac{p'\cdot\epsilon_2}{p'\cdot k_2} \Bigr)\cdots
    \Bigl( \frac{p'\cdot\epsilon_n}{p'\cdot k_n} \Bigr)
  = \bar{u}(p')\prod_{i=1}^{n}\ \Bigl( \frac{-ie\,p'\cdot\epsilon_i}
    {p'\cdot k_i} \Bigr),
  \label{eq:6.76}
\end{equation}
where the blob denotes a sum over all possible orders of inserting the $n$
photon lines.

A similar set of manipulations simplifies the sum over soft photon insertions
on the initial electron line.  There, however, the propagator momenta are
$p-k_1$, $p-k_1-k_2$, and so on.  We therefore get an extra minus sign in the
factor for each photon, since
\begin{equation}
  (p-\textstyle\sum k)^2 - m^2 \approx -2p\cdot\textstyle\sum k .
\end{equation}
Now consider diagrams containing a total of $n$ soft photons, connected in any
possible order to the initial or final electron lines.  The sum over all such
diagrams can be written
\begin{equation}
  \bar{u}(p')\,\mathcal{M}_{\mathrm{hard}}\,u(p)\ \prod_{i=1}^{n}\ (-ie)\
    \Bigl( \frac{p'\cdot\epsilon_i}{p'\cdot k_i} - \frac{p\cdot\epsilon_i}
      {p\cdot k_i} \Bigr),
  \label{eq:6.77}
\end{equation}
By multiplying out all the factors, you can see that we get the correct term
for each possible way of dividing the $n$ photons between the two lines.

Next we must decide which photons are real and which are virtual.  We can make
a virtual photon by picking two photon momenta $k_i$ and $k_j$, setting
$k_j = -k_i = k$, multiplying by the photon propagator, and integrating over
$k$.  For each virtual photon we then obtain the expression
\begin{equation}
  \begin{aligned}
    X &= \frac{1}{2}\ \frac{e^2}{2}\ \int \frac{d^4k}{(2\pi)^4}\
      \frac{-i}{k^2+i\epsilon}\ \Bigl( \frac{p'}{p'\cdot k}
        - \frac{p}{p\cdot k} \Bigr)\Bigl( \frac{p'}{-p'\cdot k}
          - \frac{p}{-p\cdot k} \Bigr)\\
    &= \frac{ie^2}{2}\ \int \frac{d^4k}{(2\pi)^4}\ \frac{1}{k^2+i\epsilon}\
      \Bigl( \frac{p'}{p'\cdot k} - \frac{p}{p\cdot k} \Bigr)^2 .
  \end{aligned}
  \label{eq:6.78}
\end{equation}
The factor of $1/2$ is required because our procedure has counted each Feynman
diagram twice: interchanging $k_i$ and $k_j$ gives back the same diagram.  It
is possible to evaluate this expression by careful contour integration, but
there is an easier way.  Notice that this approximation scheme assigns to the
diagram with one loop and no external photons the value
\begin{equation}
  \bar{u}(p')(i\mathcal{M}_{\mathrm{hard}})u(p)\cdot X .
\end{equation}
Thus, $X$ must be precisely the infrared limit of the one-loop correction to
the form factor, as displayed in \eqref{eq:6.61}:
\begin{equation}
  X = \frac{\alpha}{\pi}\mathcal{I}_{\mathrm{IR}}(q^2)\log\Bigl(
    \frac{-q^2}{\mu^2} \Bigr) .
  \label{eq:6.79}
\end{equation}
A direct derivation of this result from \eqref{eq:6.78} is given in Weinberg's
paper cited above.  Note that result \eqref{eq:6.79} followed in our argument
of the previous section only after the subtraction at $q^2 = 0$, and so we
should worry whether \eqref{eq:6.79} is consistent with the corresponding
subtraction of the $n$th-order diagram.  In addition, some of the diagrams we
are summing contain external-leg corrections, which we have not discussed.
Here we simply remark that neither of these subtleties affects the final
answer; the proof requires the heavy machinery in the paper of Yennie,
Frautschi, and Suura.

If there are $m$ virtual photons we get $m$ factors like \eqref{eq:6.79}, and
also an additional symmetry factor of $1/m!$ since interchanging virtual
photons with each other does not change the diagram.  We can then sum over
$m$ to obtain the complete correction due to the presence of arbitrarily many
soft virtual photons:
\begin{equation}
  \sum_{m=0}^{\infty}\ \frac{X^m}{m!}
    = \bar{u}(p')(i\mathcal{M}_{\mathrm{hard}})u(p)\exp(X) .
  \label{eq:6.80}
\end{equation}
If in addition to the $m$ virtual photons we also emit a real photon, we must
multiply by its polarization vector, sum over polarizations, and integrate the
squared matrix element over the photon's phase space.  This gives an
additional factor
\begin{equation}
  \int \frac{d^3k}{(2\pi)^3}\ \frac{1}{2k}\ \frac{\alpha}{2\pi}\ \sum
    \Bigl| \frac{p'\cdot\epsilon}{p'\cdot k} - \frac{p\cdot\epsilon}{p\cdot k}
      \Bigr|^2
  = \frac{\alpha}{\pi}\ \mathcal{I}(\mathbf{v},\mathbf{v}')\log\Bigl(
    \frac{E_{-}}{\mu} \Bigr)
  = Y ,
  \label{eq:6.81}
\end{equation}
in the cross section.  Assuming that the energy of the photon is greater than
$\mu$ and less than $E_{-}$ (the detector threshold), this expression is
simply
\begin{equation}
  Y = \frac{\alpha}{\pi}\mathcal{I}(\mathbf{v},\mathbf{v}')\log\Bigl(
    \frac{E_{-}}{\mu} \Bigr)
  = \frac{\alpha}{\pi}\mathcal{I}_{\mathrm{IR}}(q^2)\log\Bigl(
    \frac{E_{-}}{\mu} \Bigr) .
  \label{eq:6.82}
\end{equation}
If $n$ real photons are emitted we get $n$ such factors, and also a symmetry
factor of $1/n!$ since there are $n$ identical bosons in the final state.  The
cross section for emission of any number of soft photons is therefore
\begin{equation}
  \sum_{n=0}^{\infty}\ \frac{Y^n}{n!}\ d\sigma_{\mathrm{hard}}
    = d\sigma_{\mathrm{hard}}\ \exp(Y) .
  \label{eq:6.83}
\end{equation}
The probability of scattering without emitting a hard photon is the total
cross section, and the expression
\begin{equation}
  d\sigma_{\mathrm{total}} = d\sigma_{\mathrm{hard}}\,
    \exp(X + Y) ,
  \label{eq:6.84}
\end{equation}
is the central result of this section.  The infrared divergences of $X$ (from
virtual photons) and $Y$ (from real photons) cancel exactly, leaving a finite
result.  In the limit of large momentum transfer, the exponential factor
contains the Sudakov double logarithm, and
\begin{equation}
  \exp(X+Y) \sim \exp\Bigl[ -\frac{\alpha}{\pi}
    \mathcal{I}_{\mathrm{IR}}(q^2)\log\Bigl( \frac{-q^2}{m^2} \Bigr)
      \log\Bigl( \frac{-q^2}{E_{-}^2} \Bigr) \Bigr] .
  \label{eq:6.85}
\end{equation}
In this limit, the probability of scattering without emitting a hard photon
decreases faster than any power of $q^2$.  The exponential correction factor,
containing the Sudakov double logarithm, is known as the \emph{Sudakov form
factor}.

To conclude this section, let us calculate the probability, in the same
approximation, that some hard scattering process is accompanied by the
production of $n$ soft photons, all with energies between $E_{-}$ and $E_{+}$.
The phase-space integral for these photons gives $\log(E_{+}/E_{-})$ instead
of $\log(E_{\mathrm{max}}/\mu)$.  If we assign photons with energy greater
than $E_{+}$ to the ``hard'' part of the process, we find that the cross
section is given by \eqref{eq:6.84}, times the additional factor
\begin{equation}
  \mathrm{Prob}\bigl( n\gamma \text{ with } E_{-} < E < E_{+} \bigr)
    = \frac{1}{n!}\Bigl[ \frac{\alpha}{\pi}\mathcal{I}_{\mathrm{IR}}(q^2)
      \log\Bigl( \frac{E_{+}}{E_{-}} \Bigr) \Bigr]^n
      \exp\Bigl[ -\frac{\alpha}{\pi}\mathcal{I}_{\mathrm{IR}}(q^2)
      \log\Bigl( \frac{E_{+}}{E_{-}} \Bigr) \Bigr] .
  \label{eq:6.86}
\end{equation}
This expression has the form of a Poisson distribution,
\begin{equation}
  P(n) = \frac{1}{n!}\lambda^n e^{-\lambda},
\end{equation}
with
\begin{equation}
  \lambda = \frac{\alpha}{\pi}\mathcal{I}_{\mathrm{IR}}(q^2)\log\Bigl(
    \frac{E_{+}}{E_{-}} \Bigr) .
\end{equation}
This is precisely the semiclassical estimate of the number of radiated photons
that we made in Eq.~\eqref{eq:6.19}.

\problems

\textbf{6.1} \emph{Rosenbluth formula.}  As discussed in Section~6.2, the
exact electromagnetic interaction vertex for a Dirac fermion can be written
quite generally in terms of two form factors $F_1(q^2)$ and $F_2(q^2)$:
\begin{equation}
  -ie\Gamma^\mu(p',p) = -ie\Bigl[ \gamma^\mu F_1(q^2)
    + \frac{i\sigma^{\mu\nu}q_\nu}{2m}F_2(q^2) \Bigr],
\end{equation}
where $q = p'-p$ and $\sigma^{\mu\nu} = \tfrac{i}{2}[\gamma^\mu, \gamma^\nu]$.
If the fermion is a strongly interacting particle such as the proton, the form
factors reflect the structure that results from the strong interactions and so
are not easy to compute from first principles.  However, these form factors
can be determined experimentally.  Consider the scattering of an electron with
energy $E \gg m_e$ from a proton initially at rest.  Show that the above
expression for the vertex leads to the following expression (the
\emph{Rosenbluth formula}) for the elastic scattering cross section, computed
to leading order in $\alpha$ but to all orders in the strong interactions:
\begin{equation}
  \frac{d\sigma}{d\Omega} = \frac{\alpha^2\cos^2(\theta/2)}
    {4E^2\sin^4(\theta/2)}\ \frac{1}{1 + (2E/m_p)\sin^2(\theta/2)}\
    \Bigl[ F_1^2 - \frac{q^2}{4m_p^2}\ \Bigl( F_2^2 \Bigr)
      + \cdots \Bigr],
\end{equation}
where $\theta$ is the lab-frame scattering angle and $F_1$ and $F_2$ are to be
evaluated at the $q^2$ associated with elastic scattering at this angle.  By
measuring $(d\sigma/d\cos\theta)$ as a function of angle, it is thus possible
to extract $F_1$ and $F_2$.  Note that when $F_1 = 1$ and $F_2 = 0$, the
Rosenbluth formula reduces to the Mott formula (in the massless limit) for
scattering off a point particle (see Problem~5.1).

\textbf{6.2} \emph{Equivalent photon approximation.}  Consider the process in
which electrons of very high energy scatter from a target.  In leading order
in $\alpha$, the electron is connected to the target by one photon propagator.
If the initial and final energies of the electron are $E$ and $E'$, the photon
will carry momentum $q$ such that $q^2 \approx -2EE'(1-\cos\theta)$.  In the
limit of forward scattering, whatever the energy loss, the photon momentum
approaches $q^2 = 0$; thus the reaction is highly peaked in the forward
direction.  It is tempting to guess that, in this limit, the virtual photon
becomes a real photon.  Let us
\begin{enumerate}[label=(\alph*)]
  \item Show that, for the process in which the electron scatters off the
    target (assumed to be spinless and very heavy), the differential cross
    section can be written as the product of the photon-target cross section
    and an electron-photon flux factor.
  \item Compute the flux of virtual photons in the forward direction, and
    show that it has the form of the equivalent photon spectrum.
  \item Show that the Weizs\"acker--Williams approximation gives a good
    description of the total cross section in the limit where the photon is
    almost real.
  \item Derive the Weizs\"acker--Williams formula for the energy spectrum of
    virtual photons.
\end{enumerate}

\textbf{6.3} \emph{Exotic contributions to $g-2$.}  Any particle that couples
to the electron can produce a correction to the electron-photon form factors
and, in particular, a correction to $g-2$.  Because the electron $g-2$ agrees
with QED to high accuracy, these corrections allow us to constrain the
properties of hypothetical new particles.
\begin{enumerate}[label=(\alph*)]
  \item The unified theory of weak and electromagnetic interactions contains
    a scalar particle $h$ called the Higgs boson, which couples to the electron
    according to
    \begin{equation}
      \mathcal{L} = \lambda\bar{\psi}\psi h .
    \end{equation}
    Compute the contribution of a virtual Higgs boson to the electron
    $(g-2)$, in terms of $\lambda$ and the mass $m_h$ of the Higgs boson.
  \item QED accounts extremely well for the electron's anomalous magnetic
    moment.  If $a = (g-2)/2$,
    \begin{equation}
      |a_{\mathrm{expt.}} - a_{\mathrm{QED}}| < 1\times10^{-10} .
    \end{equation}
    What limits does this place on $\lambda$ and $m_h$?  In the simplest
    version of the electroweak theory, $\lambda = 3\times10^{-6}$ and
    $m_h > 60$ GeV.  Show that these values are not excluded.  The coupling of
    the Higgs boson to the muon is larger by a factor $(m_\mu/m_e)$:
    $\lambda = 6\times10^{-4}$.  Thus, although our experimental knowledge of
    the muon anomalous magnetic moment is not as precise,
    \begin{equation}
      |a_{\mu,\mathrm{expt.}} - a_{\mu,\mathrm{QED}}| < 3\times10^{-9},
    \end{equation}
    one can still obtain a stronger limit on $m_h$.  Is it strong enough?
  \item Some more complex versions of this theory contain a pseudoscalar
    particle called the axion, which couples to the electron according to
    \begin{equation}
      \mathcal{L} = \lambda\bar{\psi}\gamma^5\psi\,a .
    \end{equation}
    The axion may be as light as the electron, or lighter, and may couple more
    strongly than the Higgs boson.  Compute the contribution of a virtual
    axion to the $g-2$ of the electron, and work out the excluded values of
    $\lambda$ and $m_a$.
\end{enumerate}
